\section{Summary}
In this chapter, we propose \flexheapName{}, a heap resizing mechanism for Garbage First (G1), the default OpenJDK garbage collector, that dynamically partitions a fixed DRAM budget between the JVM heap and the I/O cache.
Between GC intervals, it estimates the CPU time lost to GC and to I/O stalls and repartitions DRAM to reduce their combined cost.
\flexheapName{} relies on three concepts:
%
(1) It makes resizing decisions using G1 collection boundaries.
%
(2) It uses a history-based approach to estimate the cost of \flexheapGC{} and I/O stalls for the future intervals.
%
(3) It uses an adaptive resizing step that scales with changes in the combined cost.

We implement \flexheapName{} in OpenJDK 21’s G1 garbage collector and evaluate it on two widely used systems: the Elasticsearch search engine and the Spark analytic framework.
Compared to the G1 heap resizing mechanism, \flexheapName{} improves performance by an average of 30\% in Elasticsearch and by an average of 33\% in Spark.
It outperforms Vertical G1, a state-of-the-art enhancement to the default G1 heap resizing mechanism, that returns unused memory to the OS eagerly, by 50\% on average in throughput, demonstrating that JVM heap resizing needs to consider I/O overhead in search and analytics applications.
